PC1 retains 92.5\% of standardized variance, PC2 retains 6.3\%, and their cumulative share is 98.8\%. The 80\% retention rule selects 1 component(s); the two-dimensional display is still useful for inspecting the smaller second-direction contrast.

The displayed PC1 loadings are VOA\_dB $+0.561$, QBER $+0.591$, LogSKR $-0.580$; PC2 loadings are VOA\_dB $+0.808$, QBER $-0.242$, LogSKR $+0.537$. Variables with the same loading sign move together along a component, whereas opposite signs define a contrast. The largest absolute contributions are QBER on PC1 and VOA\_dB on PC2. Component signs are arbitrary and have been oriented so each largest loading is positive; reversing an entire component changes no statistical conclusion.

Higher PC1 scores combine higher Attenuation, QBER with lower log10(SKR + 1). Higher PC2 scores combine higher Attenuation, log10(SKR + 1) with lower QBER.


PC1 can therefore be interpreted as an attenuation/error--throughput deterioration axis: it contrasts greater attenuation and error with lower key throughput. PC2 mainly captures the attenuation contrast left after that shared direction is removed. The near-one-dimensional pattern explains why the variables can be compressed, while the smaller contrasts remain relevant for distinguishing conditions.
